\documentclass[10pt]{article}
\usepackage[T1]{fontenc}
\usepackage{amsmath,amssymb,mathtools}
\usepackage[margin=0.62in]{geometry}
\pagestyle{empty}
\setlength{\parindent}{0pt}
\setlength{\parskip}{0.32em}
\setlength{\abovedisplayskip}{0.42em}
\setlength{\belowdisplayskip}{0.42em}

\newcommand{\R}{\mathbb R}
\newcommand{\Z}{\mathbb Z}
\newcommand{\dd}{\,\mathrm d}

\begin{document}

\begin{center}
{\bfseries Taylor--Green Problem from a Generic Constrained Form}
\end{center}

\textbf{Generic form.}
Let
\[
  Q:\Omega\to\R^m
\]
be the unknown field.  Here \(\Omega\) is the continuous domain on which the
field is defined, and \(m\) is the number of scalar components of the field.
Let
\[
  \mathcal Y=\{Y_1,\ldots,Y_N\}
\]
be the domains on which equations are imposed.  For each \(Y_a\), let
\[
  R_a[Q]:Y_a\to\R^{r_a},
  \qquad
  \lambda_a:Y_a\to\R^{r_a}.
\]
The integer \(r_a\) is the number of scalar equations imposed at each point of
\(Y_a\).
The constrained functional is
\[
  \boxed{
  \mathcal L(Q,\lambda)
  =
  \sum_{a=1}^N
  \int_{Y_a}
  \lambda_a(y_a)^T R_a[Q](y_a)\dd Y_a
  }
\]
where \(y_a\) denotes the coordinates on \(Y_a\).  The domains \(Y_a\), not the
grammar of the equation, distinguish an interior law, a boundary law, an
initial trace, a gauge, or an interface equation.  Variation with respect to
all multipliers gives
\[
  R_a[Q]=0\quad\text{on }Y_a,\qquad a=1,\ldots,N.
\]

\textbf{Generic finite-dimensional form.}
A discretization replaces the field domain by a finite coordinate set
\(\Omega^h\) and each residual domain by a finite coordinate set \(Y_a^h\):
\[
  Q^h:\Omega^h\to\R^m,\qquad
  Y_a^h=\{y_{a,1}^h,\ldots,y_{a,n_a}^h\}.
\]
It also gives a discrete residual and a multiplier on each discrete residual
domain:
\[
  R_a^h[Q^h]:Y_a^h\to\R^{r_a},
  \qquad
  \lambda_a^h:Y_a^h\to\R^{r_a}.
\]
With weights \(w_{a,j}\) on \(Y_a^h\), the discrete constrained functional is
\[
  \boxed{
  \mathcal L^h(Q^h,\lambda^h)
  =
  \sum_{a=1}^N
  \sum_{j=1}^{n_a}
  w_{a,j}\lambda_a^h(y_{a,j}^h)^T
  R_a^h[Q^h](y_{a,j}^h)
  }.
\]
Variation with respect to all discrete multipliers gives
\[
  R_a^h[Q^h]=0\quad\text{on }Y_a^h,\qquad a=1,\ldots,N.
\]

\textbf{Taylor--Green continuous specialization.}
Let \(\Z\) be the set of integers and set
\[
  Y_t=[0,t_f],\qquad
  Y_x=\R^2/(2\pi\Z)^2,\qquad
  \Omega=Y_t\times Y_x,\qquad
  \partial_0\Omega=\{0\}\times Y_x .
\]
For this problem,
\[
  t_f=1,\qquad \nu=0.01,\qquad
  x=(x_1,x_2)\in Y_x,\qquad
  Q=(u,v,p):\Omega\to\R^3 .
\]
The three residual domains are
\[
  Y_1=\Omega,\qquad
  Y_2=\partial_0\Omega,\qquad
  Y_3=Y_t .
\]

The law residual \(R_1[Q]:Y_1\to\R^3\) is
\[
R_1[Q]=
\begin{bmatrix}
u_t+u u_{x_1}+v u_{x_2}+p_{x_1}
  -\nu(u_{x_1x_1}+u_{x_2x_2})\\
v_t+u v_{x_1}+v v_{x_2}+p_{x_2}
  -\nu(v_{x_1x_1}+v_{x_2x_2})\\
p_{x_1x_1}+p_{x_2x_2}
  +u_{x_1}^2+2u_{x_2}v_{x_1}+v_{x_2}^2
\end{bmatrix}.
\]
The third row is part of the law; it is the pressure equation obtained from
the divergence of incompressible momentum.

The exact field used to select the solution is
\[
  Q^*=(u^*,v^*,p^*),
\]
\[
  u^*=\sin x_1\cos x_2\,e^{-2\nu t},\qquad
  v^*=-\cos x_1\sin x_2\,e^{-2\nu t},\qquad
  p^*=\frac{\cos 2x_1+\cos 2x_2}{4}e^{-4\nu t}.
\]
The initial-trace residual \(R_2[Q]:Y_2\to\R^2\) is
\[
R_2[Q]=
\begin{bmatrix}
u|_{\partial_0\Omega}-u^*|_{\partial_0\Omega}\\
v|_{\partial_0\Omega}-v^*|_{\partial_0\Omega}
\end{bmatrix}.
\]
The pressure-gauge residual \(R_3[Q]:Y_3\to\R\) is
\[
  R_3[Q](t)=\int_{Y_x}p(t,x)\dd x .
\]
Thus
\[
\mathcal L
=
\int_{Y_1}\lambda_1^T R_1[Q]\dd Y_1
+
\int_{Y_2}\lambda_2^T R_2[Q]\dd Y_2
+
\int_{Y_3}\lambda_3 R_3[Q]\dd Y_3 .
\]

\textbf{Taylor--Green finite-dimensional specialization.}
Choose discrete coordinate sets
\[
  Y_t^h=\{t_0,\ldots,t_n\},\qquad
  Y_x^h=\{x_c:c\in C_h\},
\]
and define
\[
  \Omega^h=Y_t^h\times Y_x^h,\qquad
  \partial_0\Omega^h=\{t_0\}\times Y_x^h.
\]
The discrete residual domains are
\[
  Y_1^h=\Omega^h,\qquad
  Y_2^h=\partial_0\Omega^h,\qquad
  Y_3^h=Y_t^h.
\]
The discrete field is
\[
  Q^h=(u^h,v^h,p^h):\Omega^h\to\R^3 .
\]

Let \(D_t,D_1,D_2,D_{11},D_{22}\) be the selected derivative approximations.
Write
\[
  \widehat u_t=D_tu^h,\qquad
  \widehat u_i=D_i u^h,\qquad
  \widehat u_{ii}=D_{ii}u^h,
\]
and similarly for \(v^h,p^h\).  On \(Y_1^h\),
\[
R_1^h[Q^h]=
\begin{bmatrix}
\widehat u_t+u^h\widehat u_1+v^h\widehat u_2+\widehat p_1
  -\nu(\widehat u_{11}+\widehat u_{22})\\
\widehat v_t+u^h\widehat v_1+v^h\widehat v_2+\widehat p_2
  -\nu(\widehat v_{11}+\widehat v_{22})\\
\widehat p_{11}+\widehat p_{22}
  +\widehat u_1^2+2\widehat u_2\widehat v_1+\widehat v_2^2
\end{bmatrix}.
\]
On \(Y_2^h\) and \(Y_3^h\),
\[
R_2^h[Q^h](x_c)=
\begin{bmatrix}
u^h(t_0,x_c)-u^*(t_0,x_c)\\
v^h(t_0,x_c)-v^*(t_0,x_c)
\end{bmatrix},
\qquad
R_3^h[Q^h](t_k)=\sum_{c\in C_h}w_c\,p^h(t_k,x_c).
\]

With quadrature weights \(w_{a,j}\) on each \(Y_a^h\), the discrete
functional is
\[
  \boxed{
  \mathcal L^h(Q^h,\lambda^h)
  =
  \sum_{a=1}^3
  \sum_{y_{a,j}^h\in Y_a^h}
  w_{a,j}\lambda_a^h(y_{a,j}^h)^T
  R_a^h[Q^h](y_{a,j}^h)
  }.
\]
Variation with respect to all discrete multipliers gives
\[
  R_a^h[Q^h]=0\quad\text{on }Y_a^h,\qquad a=1,2,3.
\]


\section{The jet along the design}

A \emph{jet} is the list of the derivatives of a quantity along a
coordinate at one point, to some order: the coefficients of its
Taylor expansion there, the expansion not summed. For the unknown
$u(t, \nu)$ at a fixed instant $t_k$ and the design value $\nu_0$, the
jet along $\nu$ of order three is
\[
  \Bigl( u,\; \frac{\partial u}{\partial \nu},\;
         \frac{\partial^2 u}{\partial \nu^2},\;
         \frac{\partial^3 u}{\partial \nu^3} \Bigr)(t_k, \nu_0),
\]
and it represents $u(t_k, \nu)$ near $\nu_0$ through
\[
  u(t_k, \nu) = u + \frac{\partial u}{\partial \nu}(\nu - \nu_0)
  + \frac{\partial^2 u}{\partial \nu^2}\frac{(\nu - \nu_0)^2}{2}
  + \frac{\partial^3 u}{\partial \nu^3}\frac{(\nu - \nu_0)^3}{6} + \cdots
\]

The word is the one used along time. The discrete residual stores, at
every instant, the jet of $u$ along $t$, $(u, u_t, u_{tt})$, which the
time families connect. The design coordinate $\nu$ is a point of the
design space, fixed at $\nu_0$; its discretization is the expansion of
an order at that value, so that the discrete manifold
$\Omega_h = \{t_1, \ldots, t_n\} \times \mathrm{expansion}(3)$ is the
grid
\[
\begin{array}{c|cccc}
       & \text{order } 0 & \text{order } 1 & \text{order } 2 & \text{order } 3 \\
       & \text{(the state)} & \dfrac{d}{d\nu} & \dfrac{d^2}{d\nu^2} & \dfrac{d^3}{d\nu^3} \\[1ex]
\hline
  t_n  & [u, u_t, u_{tt}] & [u, u_t, u_{tt}]' & [u, u_t, u_{tt}]'' & [u, u_t, u_{tt}]''' \\
  \vdots & \vdots & \vdots & \vdots & \vdots \\
  t_2  & [u, u_t, u_{tt}] & [u, u_t, u_{tt}]' & [u, u_t, u_{tt}]'' & [u, u_t, u_{tt}]''' \\
  t_1  & [u, u_t, u_{tt}] & [u, u_t, u_{tt}]' & [u, u_t, u_{tt}]'' & [u, u_t, u_{tt}]'''
\end{array}
\]
Column $0$ is the solution. Column $1$ is $du/d\nu$ at every instant,
column $2$ is $d^2u/d\nu^2$, and so on. The derivative of a discrete
field along $\nu$ is the field of the next column, with the columns
after it as its jet; in the programs, \texttt{u\_h \% derivative([nu])}.

The multipliers have the same columns. The adjoint $\lambda_r(t)$ at
every instant has its jet along $\nu$ beside it, and $\kappa$, the
multiplier of the condition $\nu - \nu_0 = 0$, a single number, has a
jet of four numbers. Since $\kappa = -dJ/d\nu$, that jet is
\[
  \Bigl( -\frac{dJ}{d\nu},\; -\frac{d^2J}{d\nu^2},\;
         -\frac{d^3J}{d\nu^3},\; -\frac{d^4J}{d\nu^4} \Bigr),
\]
the sensitivities of the objective by the adjoint, which the expansion
of the state computes independently.

The columns follow from the equations holding at every $\nu$. Let
$F(u(\nu), \nu) = 0$ be the discrete residual over a block and
$A = \partial F / \partial u$ its Jacobian. Differentiating once along
$\nu$,
\[
  A\, u' + \frac{\partial F}{\partial \nu} = 0,
\]
so column $1$ is one linear solve with $A$. Differentiating twice,
\[
  A\, u'' + \frac{d}{d\nu}\Bigl(\frac{\partial F}{\partial u}\Bigr) u'
          + \frac{d}{d\nu}\Bigl(\frac{\partial F}{\partial \nu}\Bigr) = 0,
\]
so column $2$ is one more solve with the same $A$, its right side in
$u$ and $u'$; and so on to the order. The jet arithmetic of the type
\texttt{derivative\_terms} evaluates $F$ on the jets of $u$ and $\nu$,
with the coefficient of order $m$ of $u$ set to zero, and returns the
jet of $F$, whose coefficient at order $m$ is exactly the right side of
the $m$-th solve. Nothing is differentiated by hand and nothing is
approximated by differences.

The adjoint's columns follow the same way from the adjoint equation
$A^{\mathsf T} \mu = -\partial J / \partial u$: differentiating along
$\nu$ leaves $A^{\mathsf T}$ on the left, so each order is one more
transposed solve per block, its right side formed from the jets of the
state and the lower columns of $\mu$.

\end{document}
